\documentclass[11pt,a4paper]{article}
\usepackage[utf8]{inputenc}
\usepackage[T1]{fontenc}
\usepackage[french]{babel}
\usepackage{lmodern}
\usepackage{amsmath,amssymb}
\usepackage{geometry}
\usepackage{booktabs}
\usepackage{hyperref}
\usepackage{xcolor}
\usepackage{enumitem}
\usepackage{microtype}
\usepackage{tcolorbox}

\geometry{margin=2.4cm}
\hypersetup{
  colorlinks=true,
  linkcolor=blue!55!black,
  urlcolor=blue!55!black,
  citecolor=blue!55!black
}
\setlength{\parindent}{0pt}
\setlength{\parskip}{0.6em}
\setlist[itemize]{leftmargin=1.4em}

\newcommand{\lp}{\ell_{\rm P}}
\newcommand{\rhoP}{\rho_{\rm P}}
\newcommand{\rhoDE}{\rho_{\rm DE}}
\newcommand{\rhoBH}{\rho_{\rm BH}}
\newcommand{\RL}{R_\Lambda}

\title{Programme de tests experimentaux du cadre GUP--UV/IR\\
\large Beta0, holographie, choix de l'horizon, energie noire dynamique et relation alpha--Lambda}
\author{Document de travail}
\date{22 mai 2026}

\begin{document}
\maketitle

\begin{abstract}
Cette note transforme les consequences du cadre GUP--UV/IR en programme de
tests. Les resultats deja obtenus donnent une conclusion nette : un parametre
\(\beta_0\sim10^{60}\) ajuste directement a l'energie noire est exclu comme
parametre local universel. La branche viable est une architecture en deux
etapes : regularisation UV locale puis projection gravitationnelle ou
holographique IR. Le choix de la longueur IR est fortement contraint par la
non-variation observee de \(\alpha\), ce qui defavorise \(L=c/H(t)\) et rend
plus naturelle une echelle quasi constante comme \(R_\Lambda\).
\end{abstract}

\tableofcontents
\newpage

\section{Test 1 : exclure un \texorpdfstring{\(\beta_0\sim10^{60}\)}{beta0~10^60} universel}

La regularisation GUP du vide massless donne
\begin{equation}
\rho_{\rm reg}=\frac{g\rhoP}{16\pi^2\beta_0^2}.
\end{equation}

Identifier directement \(\rho_{\rm reg}=\rhoDE\) donne, pour \(g=1\),
\begin{equation}
\beta_{0,\rm DE}\simeq2.36\times10^{60}.
\end{equation}

Cette valeur est exclue si elle est universelle. Les tests publics deja
effectues donnent :
\begin{center}
\begin{tabular}{lcc}
\toprule
Canal & Borne indicative sur \(\beta_0\) & Verdict \\
\midrule
FIRAS/CMB & \(3.7\times10^{58}\) & exclut \(\beta_{0,\rm DE}\) \\
BBN/radiation & \(2.6\times10^{41}\) & exclut fortement \\
Fermions denses & \(1.7\times10^{38}\) & exclut fortement \\
Electron atomique & \(\chi_e/\beta_0=9.33\times10^{-50}\) & exclut \(\beta_{0,\rm DE}\) \\
\bottomrule
\end{tabular}
\end{center}

Pour l'electron atomique,
\begin{equation}
\chi_e=\beta_0\left(\frac{m_e}{m_P}\right)^2\alpha^2.
\end{equation}
Donc \(\beta_{0,\rm DE}\) donnerait \(\chi_e\sim2.2\times10^{11}\), ce qui est
incompatible avec la physique atomique.

\begin{tcolorbox}[colback=red!3!white,colframe=red!50!black,title=Conclusion]
\(\beta_0\sim10^{60}\) ne peut pas etre un parametre GUP local universel. Il ne
peut etre qu'un parametre effectif, sectoriel ou holographique.
\end{tcolorbox}

\section{Test 2 : branche holographique}

La borne gravitationnelle d'une region de taille \(L\) donne
\begin{equation}
\boxed{
\rhoBH(L)=\frac{3c^4}{8\pi G L^2}.
}
\end{equation}

Pour \(L_H=c/H_0\),
\begin{equation}
\rhoBH(L_H)=\frac{3c^2H_0^2}{8\pi G}=\rho_c.
\end{equation}

Cette relation est robuste, mais elle est seulement une borne ou une identite
dimensionnelle. Le probleme dynamique reste :
\begin{equation}
\rhoDE=\Omega_\Lambda\rho_c,
\qquad
w\simeq -1.
\end{equation}

Il faut donc expliquer pourquoi la partie gravitationnellement active sature
une fraction \(\Omega_\Lambda\) de la borne et pourquoi le tenseur effectif a
une pression negative.

\section{Test 3 : choix de la longueur IR}

La relation alpha--IR s'ecrit
\begin{equation}
\alpha^{-1}(L)=\ln\left(\frac{L}{10\pi\lp}\right).
\end{equation}

Si \(L\) varie, alors
\begin{equation}
\frac{\dot\alpha}{\alpha}
=-\alpha\,\frac{\dot L}{L}.
\end{equation}

Pour le choix naif \(L=c/H(t)\), en cosmologie plate \(\Lambda\)CDM,
\begin{equation}
\frac{\dot L}{L}=-\frac{\dot H}{H}
\simeq\frac32\Omega_m H_0.
\end{equation}

Avec les valeurs Planck 2018, cela donne
\begin{equation}
\left|\frac{\dot\alpha}{\alpha}\right|
\simeq2.4\times10^{-13}\ {\rm an}^{-1}.
\end{equation}

Les horloges atomiques contraignent typiquement cette variation autour de
\(10^{-18}\ {\rm an}^{-1}\). Donc \(L=c/H(t)\) est fortement defavorise.

\begin{tcolorbox}[colback=yellow!4!white,colframe=yellow!50!black,title=Conclusion]
La longueur IR pertinente pour la relation \(\alpha\)--\(\Lambda\) doit etre
quasi constante : \(R_\Lambda\), horizon de de Sitter asymptotique, ou horizon
futur avec variation tres faible.
\end{tcolorbox}

\section{Test 4 : energie noire dynamique}

Un modele holographique standard donne souvent
\begin{equation}
w=-\frac13-\frac{2}{3c_{\rm hde}}\sqrt{\Omega_{\rm DE}}.
\end{equation}

Pour \(c_{\rm hde}=1\) et \(\Omega_{\rm DE}\simeq0.685\),
\begin{equation}
w\simeq -0.885.
\end{equation}

Cette valeur est proche de \(-1\), mais deja tendue face aux contraintes
cosmologiques modernes. Le test naturel est donc :
\begin{equation}
w(a)\neq -1
\quad\text{mais}\quad
w(a)\approx -1.
\end{equation}

Les donnees BAO, supernovae, CMB et croissance des structures peuvent tester
cette branche.

\section{Test 5 : relation \texorpdfstring{\(\alpha\)--\(\Lambda\)}{alpha--Lambda}}

La relation candidate est
\begin{equation}
\boxed{
\Lambda=
\frac{3}{100\pi^2\lp^2}e^{-2/\alpha}.
}
\end{equation}

Avec CODATA 2022 et Planck 2018, on obtient
\begin{equation}
\alpha^{-1}_{\rm pred}=137.036063742708,
\qquad
\alpha^{-1}_{\rm obs}=137.035999177,
\end{equation}
et
\begin{equation}
\frac{\Lambda(\alpha)}{\Lambda_{\rm obs}}=1.000129139753.
\end{equation}

La coherence numerique est forte. Elle devient predictive seulement si le
facteur \(10\pi\), ou \(g_{\rm eff}^{U(1)}=5/2\), est derive
independamment.

\section{Test 6 : mode de bord Maxwell}

Le verrou theorique est
\begin{equation}
\boxed{
g_{\rm edge}^{U(1)}=\frac12.
}
\end{equation}

Ce n'est pas un test de donnees publiques. C'est un calcul de theorie quantique
des champs avec bord. La structure attendue est
\begin{equation}
Z_{\rm edge}\propto
\left[\det{}'(-\Delta_{S^2})\right]^{\pm1/2}.
\end{equation}

Il faut montrer que la capacite holographique renormalisee du canal de bord
vaut bien \(1/2\), de maniere independante de la valeur observee de
\(\alpha\).

\section{Synthese}

\begin{center}
\begin{tabular}{p{0.3\linewidth} p{0.58\linewidth}}
\toprule
Question & Verdict \\
\midrule
\(\beta_0\sim10^{60}\) universel & Exclu par FIRAS, BBN, fermions denses et atomes \\
\(\rhoBH(L)\propto L^{-2}\) & Correct comme borne, dynamique ouverte \\
\(L=c/H(t)\) & Fortement defavorise par \(\dot\alpha/\alpha\) \\
\(w(a)\approx -1\) & Testable avec BAO/SNe/CMB \\
\(\alpha\)--\(\Lambda\) & Accord numerique fort, preuve manquante du \(10\pi\) \\
Mode de bord Maxwell & Calcul theorique central a faire \\
\bottomrule
\end{tabular}
\end{center}

\section{Conclusion}

Le programme experimental clarifie le statut du cadre. La piste directe
\(\rho_{\rm reg}=\rhoDE\) avec un \(\beta_0\) universel est exclue. La piste
viable est une architecture UV/IR : le GUP regularise l'UV, puis une projection
gravitationnelle ou holographique selectionne la partie IR active.

La contrainte nouvelle la plus forte est la variation de \(\alpha\). Elle
indique que la longueur IR de la relation \(\alpha\)--\(\Lambda\) ne peut pas
etre simplement \(c/H(t)\). Une echelle quasi constante comme \(R_\Lambda\) est
beaucoup plus naturelle.

\begin{thebibliography}{9}
\bibitem{CODATA2022}
NIST/CODATA, \emph{CODATA recommended values of the fundamental physical
constants: 2022}.
\bibitem{Planck2018}
Planck Collaboration, \emph{Planck 2018 results. VI. Cosmological parameters},
arXiv:1807.06209.
\bibitem{CKN}
A. Cohen, D. Kaplan, A. Nelson, \emph{Effective field theory, black holes, and
the cosmological constant}, Phys. Rev. Lett. 82, 4971 (1999).
\bibitem{LiHDE}
M. Li, \emph{A model of holographic dark energy}, Phys. Lett. B 603, 1 (2004).
\bibitem{DonnellyWall}
W. Donnelly and A. C. Wall, \emph{Geometric entropy and edge modes of the
electromagnetic field}, arXiv:1506.05792.
\end{thebibliography}

\end{document}
